% Consolidated guide-selftest06 from 97b62cc61307eba158a4498b55de3a5fdef3ff43; source fragments preserved.
% 2026-11-24 八批严格二校第7批：历史片段哈希对应修订前版本。

% === BEGIN guide-selftest06-choices.tex blob 8057eb62f9cfe48180174d764fbe96e6cfbbf5c0 ===
% 第17批（自测题6）：选择6.1.1—6.1.10；原书印刷163—164/PDF171—172。
\subsection{自测题6：电磁场基础选择题}

\sourcetest{6.1.1}{163}{171}
\problemtext $x$ 轴上左半无限长线电荷密度 $+\lambda$、右半无限长密度 $-\lambda$，求 $P(0,a)$ 的电场。
\workings 由库仑场沿两段线分别积分，竖直分量抵消，而水平分量同向。每段贡献 $E_x=\lambda/(4\pi\varepsilon_0a)$，故
\[\boxed{\bm E(0,a)=\frac{\lambda}{2\pi\varepsilon_0a}\,\bm i}.\]
原选择题选 \textbf{B}。
\crosscheck 正电左段排斥、负电右段吸引，二者都指向 $+x$；若只把两段净线电荷“相加为0”，会错误地得出零场。

\sourcetest{6.1.2}{163}{171}
\problemtext 正点电荷 $+q$ 到 $P$ 的距离为 $a$、到 $M$ 的距离为 $2a$，把 $P$ 选作电势零点，求 $M$ 电势。
\workings 以无穷远为零时 $\varphi(r)=q/(4\pi\varepsilon_0r)$；变更电势参考值要整体减去 $\varphi(P)$，故
\[\boxed{\varphi_M=\frac{q}{4\pi\varepsilon_0}\left(\frac1{2a}-\frac1a\right)=-\frac{q}{8\pi\varepsilon_0 a}}.\]
选 \textbf{D}。
\crosscheck 电势差不依参考零点；若直接写正值 $q/(8\pi\varepsilon_0a)$，忽略了题目重新设定的零势点。

\sourcetest{6.1.3}{163}{171}
\problemtext 同心薄球壳，内壳半径 $R_1$ 带电 $Q$，外壳半径 $R_2$ 最初中性但接地；求两壳间 $R_1<r<R_2$ 的场强与电势。
\workings 对高斯球，$\bm E=\frac{Q}{4\pi\varepsilon_0r^2}\hat{\bm r}$。外壳接地意味着 $V(R_2)=0$，对径向电场积分：
\[\boxed{V(r)=\frac{Q}{4\pi\varepsilon_0}\left(\frac1r-\frac1{R_2}\right)},\quad \boxed{E(r)=\frac{|Q|}{4\pi\varepsilon_0r^2}}.\]
与选项 \textbf{C} 相符。
\crosscheck 电势参考是接地外壳，不是无穷远；场强大小为正，方向由 $Q$ 符号决定。

\sourcetest{6.1.4}{163}{171}
\problemtext 等边三角形边长 $a$，三顶点分别放电荷 $q,2q,3q$，将另一电荷 $Q$ 从无穷远缓慢移到中心，求外力功。
\workings 三顶点到重心的距离同为 $a/\sqrt3$，因此中心电势 $V_O=\frac{1}{4\pi\varepsilon_0}\frac{(q+2q+3q)\sqrt3}{a}$。准静态移入外力功 $\Delta U=QV_O$，即
\[\boxed{W_{\rm ext}=\frac{6\sqrt3\,qQ}{4\pi\varepsilon_0a}}.\]
选 \textbf{C}。
\crosscheck 应当相加的是标量电势，不是中心处三个矢量电场。

\sourcetest{6.1.5}{163}{171}
\problemtext 质子在匀强磁场 $B=0.3$T 中垂直于磁场作半径 $R=10$cm 的圆周运动，估计动能数量级（质量 $1.67\times10^{-27}$kg，电荷 $1.6\times10^{-19}$C）。
\workings 洛伦兹力充当向心力：$qvB=mv^2/R$，故 $p=mv=qBR$、$K=q^2B^2R^2/(2m)$。代入
\[\boxed{K\approx6.90\times10^{-15}\ {\rm J}\approx0.0431\ {\rm MeV}}.\]
按十进数量级（对数距离最小，$0.0431>\sqrt{0.01\times0.1}$ MeV）选\textbf{B}（$0.1$ MeV）；若题意为线性绝对误差最小，则$0.01$ MeV更近，原选项措辞须核。
\crosscheck 题问的是数量级，而不是精确数值；$v=qBR/m\simeq2.87\times10^6$ m/s，使用非相对论动能成立。

\sourcetest{6.1.6}{164}{172}
\problemtext $N$ 匝边长为 $a$ 的正三角形通电线圈位于匀强磁场中，线圈法线与 $\bm B$ 同向，求磁力矩大小。
\workings 任意平面线圈磁矩 $\bm m=NI S\,\hat{\bm n}$，其中 $S=\sqrt3 a^2/4$。磁力矩 $\bm\tau=\bm m\times\bm B$，因 $\hat{\bm n}\parallel\bm B$，
\[\boxed{|\bm\tau|=NISB\sin0=0}.\]
选 \textbf{D}。
\crosscheck 有磁通和磁矩不代表有磁力矩；仅当磁矩与磁场既不平行也不反平行时才有非零转矩。

\sourcetest{6.1.7}{164}{172}
\problemtext 无限长薄平电流带宽度 $a$，总电流 $I$ 沿带长方向均匀分布，带外离近边缘 $b$ 的点 $P$ 磁场多大？
\workings 设$P$与薄带共面且在宽度方向的外侧，$b>0$。把电流带分解成宽 $dx$ 的无限长线电流 $dI=(I/a)\,dx$。同侧场方向一致，取从近边界量起的坐标 $x\in[0,a]$：
\[B(P)=\frac{\mu_0 I}{2\pi a}\int_0^a\frac{dx}{b+x}
=\boxed{\frac{\mu_0 I}{2\pi a}\ln\frac{a+b}{b}}.\]
选 \textbf{B}。
\crosscheck 必须按每条细线到 $P$ 的距离 $b+x$ 积分；直接把全部电流集中在中心只给近似。

\sourcetest{6.1.8}{164}{172}
\problemtext 原一匝圆线圈半径 $R$、电流 $I$，将原导线在总长度不变下绕成两匝同向圆线圈，并通以同样的 $I$；求中心磁场与磁矩变化倍数。
\workings 忽略引线长度、匝间距及线径，新两匝同心且近似共面。原长 $2\pi R$，新两匝各半径 $R'=R/2$。中心磁场 $B'=\mu_0NI/(2R')=\mu_0(2I)/R=4B$；磁矩大小 $m'=NI\pi R'^2=2I\pi R^2/4=m/2$。故
\[\boxed{\frac{B'}B=4,\qquad \frac{m'}m=\frac12}.\]
选 \textbf{B}。
\crosscheck 中心磁场对匝数和半径均敏感；磁矩与匝数乘面积成正比。

\sourcetest{6.1.9}{164}{172}
\problemtext 导体棒 $AB$ 绕棒上的 $C$ 点、以垂直棒并与匀强磁场平行的轴匀速旋转；$BC=L/3$，比较端点 $A,B$ 的电势。
\workings 取 $\bm\omega$ 与 $\bm B$ 同向，棒上电荷感受 $\bm v\times\bm B$ 沿远离旋转轴的径向。端点相对轴的距离 $CA=2L/3$、$CB=L/3$，所以
\[\boxed{V_A-V_B=\tfrac12B\omega(CA^2-CB^2)=\frac{B\omega L^2}{6}>0}.\]
按原图转向选 \textbf{A}。
\crosscheck 电势并非单由端点速度线性比较，应按 $\omega Br\,dr$ 积分；若反转角速度与磁场相对方向，电势高低反转。

\sourcetest{6.1.10}{164}{172}
\problemtext 导体杆在垂直纸面向里的匀强磁场中沿轨道向右匀加速运动，外接铁芯绕组和电容器；判断达到稳定运动阶段后电容器 $M$ 极板的电荷符号和随时间变化。
\workings 磁通增加时感应电流方向由楞次定律确定：回路倾向产生纸面向外的磁场。若电容器经铁芯耦合，其极板电势要继续使用绕组的实际接线次序和感应电动势的符号判定。只根据杆运动方向可确认其来源电动势 $|\mathcal E|=B\ell v$ 随 $v$ 增大，却\emph{不能}不看绕组极性而确定 $M$ 正负。若原边是线性$R_p$、$L_p$回路、杆匀加速$a$，则$L_p\dot I_p+R_p I_p=B\ell(v_0+at)$（选一致正向），暂态消失后$\dot I_p\to B\ell a/R_p$（$R_p>0$）。次边开路感应电压$-M_{sp}\dot I_p$趋常量，而非随原边电动势持续线性增大。若次边电容经有限电阻连接且暂态衰减，则其最终电荷趋$C$乘该有向电压；若存在理想无耗振荡或铁芯饱和，则另行建模。这里保留 \textbf{源图接线二校}，不擅定 A/B/C/D。
\crosscheck 必须逐匝确认绕组绕向；电容器的极性和是否趋恒定取决于绕向、实际电路拓扑、阻抗和磁芯线性条件，不能把理想磁耦合直接等同导线相接。
% === END guide-selftest06-choices.tex ===

% === BEGIN guide-selftest06-blanks.tex blob 49ddb4f9583035f7d75a085616e6579ff7ad8cfd ===
% 第17批：自测题6填空6.2.1—6.2.9；印刷165页／PDF173页。
\subsection{自测题6：静电场、磁场与麦克斯韦方程填空}

\sourcetest{6.2.1}{165}{173}
\problemtext 均匀带电直线段长 $d$、电荷线密度 $+\lambda$，以线段中点为球心作半径 $R>d$ 的球面，求通量和线段延长线上球面交点 $P$ 的场强。
\workings 高斯定律给出封闭球面的电通量
\[\boxed{\Phi_E=\frac{\lambda d}{\varepsilon_0}}.\]
以线段中心为 $x=0$，$P=(-R,0)$，由沿 $x$ 的库仑积分，
\[\boxed{|\bm E(P)|=\frac{\lambda}{4\pi\varepsilon_0}
\left(\frac1{R-d/2}-\frac1{R+d/2}\right)
=\frac{\lambda d}{4\pi\varepsilon_0(R^2-d^2/4)}}.\]
正电荷时 $\bm E(P)$ 沿负 $x$ 方向。
\crosscheck 球面通量只取决于包围总电荷，但局部场强不能把线段直接替换成中心点电荷；要求 $R>d/2$，题设 $R>d$ 足够。

\sourcetest{6.2.2}{165}{173}
\problemtext 两无限大平行带电面 $A,B$，两面之间场强 $E_0$ 向左；左外侧 $E_0/3$ 向左，右外侧 $E_0/3$ 向右。求两面面电荷密度。
\workings 令右向为正。跨越电荷面场法向跳跃满足 $E_{\rm right}-E_{\rm left}=\sigma/\varepsilon_0$。面 $A$：$\sigma_A=\varepsilon_0(-E_0+E_0/3)$；面 $B$：$\sigma_B=\varepsilon_0(E_0/3+E_0)$，故
\[\boxed{\sigma_A=-\frac23\varepsilon_0E_0,\qquad \sigma_B=+\frac43\varepsilon_0E_0}.\]
\crosscheck 带电片不是必然等量异号；外侧场非零意味着总面电荷密度和不为零。

\sourcetest{6.2.3}{165}{173}
\problemtext 负电荷由无穷远靠近一块不带净电荷、与地绝缘的导体，问导体内部电场以及导体电势怎样变化。
\workings 在没有其他外部源、孤立有限导体、以无穷远为零的条件下，静电平衡时导体材料内部 $\boxed{\bm E=0}$，故场强保持零。外部负电荷的靠近，使孤立中性导体整体感应到较低电势（由静电互易定理$CU=q_{\rm ext}u(\bm r)$，其中$u$为同一导体置于单位电势时的无源外部调和电势，$C>0$、$u>0$，故$U<0$；沿任意复杂路径未必单调，只有相对于远处初态必降低），故 \textbf{电势减小}。
\crosscheck 导体即使净电荷保持零，也会发生表面电荷重新分布；“孤立”与“接地”边界条件不同。

\sourcetest{6.2.4}{165}{173}
\problemtext 两平行无限长直导线 $A,B$ 电流大小相同、方向均垂直纸面向外、间距 $a$，$P$ 为二者中点。另有图示仅包围 $A$ 的闭合回路 $l$，求 $P$ 处磁感应强度及闭路积分。
\workings 在中点，两条直线对 $P$ 的磁场方向相反、模相等，故 $\boxed{\bm B(P)=0}$。安培环路定律以面法向朝纸面外为正，图中回路只围住 $A$：
\[\boxed{\oint_l\bm B\cdot d\bm l=\mu_0 I}\]
（若取反向绕行则为负）。
\crosscheck 不能由一点 $\bm B(P)=0$ 断言整条回路积分为零；必须数穿过回路围成面的总电流。

\sourcetest{6.2.5}{165}{173}
\problemtext 平均周长0.50m 的闭合细铁环绕有500匝线圈，电流0.30A，铁芯相对磁导率600，求环中 $H,B$。
\workings 环路定律给出 $\oint\bm H\cdot d\bm l=NI$；在细环近似下
\[\boxed{H=\frac{NI}{\ell}=\frac{500(0.3)}{0.50}=300\ {\rm A/m}}.\]
且 $B=\mu_0\mu_rH$，故
\[\boxed{B\simeq(4\pi\times10^{-7})(600)(300)=0.226\ {\rm T}}.\]
\crosscheck 铁磁性材料的相对磁导率可能非线性变化，此题按题给常数 $\mu_r$ 的线性近似。

\sourcetest{6.2.6}{165}{173}
\problemtext 条形磁铁插入与冲击电流计串联的金属环，经过仪器的总感应电荷量 $q=2.0\times10^{-5}$C，回路总电阻25Ω，求磁通变化量。
\workings 对单匝恒定电阻回路，若自感不可忽略，写$RI+L\dot I=-\dot\Phi_{\rm ext}$。在初末电流均为零（或相等）的完整暂态测量中，自感项积分为零；忽略自感时直接有$I=-\dot\Phi_{\rm ext}/R$。积分得 $q=\int I\,dt=-\Delta\Phi/R$，所以
\[\boxed{|\Delta\Phi|=Rq=25(2.0\times10^{-5})=5.0\times10^{-4}\ {\rm Wb}}.\]
\crosscheck 这里$q$须为带符号积分的绝对值，若电流中途反向，不等于无符号累计电量，因此结果报告磁通变化量的\emph{大小}；有向符号由线圈取向定义。

\sourcetest{6.2.7}{165}{173}
\problemtext 一根向上通电的无限长直导线右侧，水平金属棒 $AB$ 与导线共面，且沿导线方向向上以速度 $v$ 移动。比较两端感应电势，并讨论电流反向或棒改为平行导线的情况。
\workings 原电流向上时，右侧磁场指入纸面，$\bm v\times\bm B$ 向左，故 \textbf{$V_A>V_B$}；电流反向则 \textbf{$V_B>V_A$}。若棒沿导线方向放置，$\bm v\times\bm B$ 仍水平且与棒正交，因此 $\boxed{V_A=V_B}$（仅指运动电动势贡献）。
\crosscheck 这里比较的是运动棒中电荷受洛伦兹磁力分离产生的电势差；同时假定导体无外加回路电压。

\sourcetest{6.2.8}{165}{173}
\problemtext 无限长直导线载电流 $I$，求距轴线 $a$ 处真空磁场的能量密度。
\workings $B(a)=\mu_0I/(2\pi a)$，真空磁场能量密度
\[\boxed{w_m=\frac{B^2}{2\mu_0}=\frac{\mu_0I^2}{8\pi^2a^2}}.\]
\crosscheck 该式是局部能量密度；若求整条无限长导线外部总能量，还需要外半径截断，否则发散。

\sourcetest{6.2.9}{165}{173}
\problemtext 给定麦克斯韦积分式①电场高斯定律、②法拉第定律、③磁场高斯定律、④安培—麦克斯韦定律；判断“变化的磁场伴随电场”“磁感线没有起止”“电荷伴随电场”各对应哪式。
\workings 由 $\oint\bm E\cdot d\bm l=-d\Phi_B/dt$ 得变化磁场伴随涡旋电场，选 \textbf{②}；$\oint\bm B\cdot d\bm S=0$ 表明磁感线无始无终，选 \textbf{③}；$\oint\bm D\cdot d\bm S=Q_{\rm free,enclosed}$ 表明自由电荷是电位移散度源；真空中等价于电场高斯定律，选 \textbf{①}。
\[\boxed{(1)\ \text{②};\quad(2)\ \text{③};\quad(3)\ \text{①}}\]
\crosscheck 方程①中的$\bm D$版本仅计自由电荷，$\bm E$版本计总电荷；有电荷必有电场结构，并不意味电荷附近\emph{每一点}的电场都非零；安培—麦克斯韦定律④反映电流和变化电场产生磁场。
% === END guide-selftest06-blanks.tex ===

% === BEGIN guide-selftest06-solutions.tex blob dff66fb86cacaa730db47471f3d97ee7103ef012 ===
% 第17批：自测题6计算6.3.1—6.3.3与证明6.4.1—6.4.2；印刷166/PDF174。
\subsection{自测题6：电场积分、动生电动势与安培力证明}

\sourcetest{6.3.1}{166}{174}
\problemtext 半径为 $R$ 的无限长均匀带电半圆柱面，沿轴线每单位长度上的总电荷为 $\lambda$，求轴线上场强。
\workings 设带电半圆柱的弧角 $\theta\in[-\pi/2,\pi/2]$，单位轴向长度、弧角宽 $d\theta$ 内线电荷 $d\lambda=(\lambda/\pi)d\theta$。在轴线上视各窄带为无限长线电荷，其场强大小为 $dE=d\lambda/(2\pi\varepsilon_0 R)$，方向背离该带，关于半圆对称轴的横向分量相加、垂直分量抵消。因此
\[
E=\frac{\lambda}{2\pi^2\varepsilon_0R}
\int_{-\pi/2}^{\pi/2}\cos\theta\,d\theta
=\boxed{\frac{\lambda}{\pi^2\varepsilon_0R}}
\]
（取 $\lambda>0$ 时为朝向没有带电柱面的一侧）。
\crosscheck 不能直接套用完整均匀圆柱面轴线场强为零的结论；半圆对称性只会消掉垂直于对称轴的分量。

\sourcetest{6.3.2}{166}{174}
\problemtext 半径 $r_1=10$cm、$r_2=20$cm 的两同心球面具有相同面电荷密度 $\sigma$，取无穷远电势零，球心电势为300V。求 $\sigma$；再求使球心电势降为零须从外球面去掉多少净电荷。
\workings 各均匀带电薄球壳内部电势恒为壳表面值。故
\[
U_O=\frac{1}{4\pi\varepsilon_0}
\left(\frac{4\pi r_1^2\sigma}{r_1}+
\frac{4\pi r_2^2\sigma}{r_2}\right)
=\frac{\sigma(r_1+r_2)}{\varepsilon_0}.
\]
代入数据
\[\boxed{\sigma=\frac{\varepsilon_0(300)}{0.30}
=8.85\times10^{-9}\ {\rm C/m^2}}.\]
若内球面电量固定、外球面电荷变化仍保持球对称（或用球心处距所有外球面电荷相同的性质），外球面改变 $\Delta Q_2$，球心电势变 $\Delta U=\Delta Q_2/(4\pi\varepsilon_0r_2)$。要求 $\Delta U=-300$V，所以需去掉\emph{有符号的正电荷量}
\[
\boxed{Q_{\rm removed}=4\pi\varepsilon_0r_2(300)
\simeq6.67\times10^{-9}\ {\rm C}}.
\]
\crosscheck 所去正电荷量大于原外球面初始正电量，因此最后外壳需带净负电（相当于继续向外壳补入电子）；“放掉原有正电荷而不允许变负”则不能实现零电势，须交代此物理限制。

\sourcetest{6.3.3}{166}{174}
\problemtext 竖直无限长直导线载上行电流 $I$；其右侧同平面内半径为 $b$ 的半圆形导电弧 $MN$，圆心距直导线 $a>b$，整段弧以速度 $v$ 向上移动。求电动势及端点电压 $V_M-V_N$。
\workings 右侧磁场向纸面内且随离导线距离 $x$ 变化：$B(x)=\mu_0I/(2\pi x)$。移动导体的洛伦兹动生电场 $\bm v\times\bm B$ 向左，沿弧从左端 $M$ 到右端 $N$ 的电动势
\[
\mathcal E_{M\to N}=\int_M^N(\bm v\times\bm B)\cdot d\bm l
=-\int_{a-b}^{a+b}\frac{\mu_0Iv}{2\pi x}\,dx
=\boxed{-\frac{\mu_0Iv}{2\pi}\ln\frac{a+b}{a-b}}.
\]
开路时左端积累正电荷，
\[
\boxed{V_M-V_N=\frac{\mu_0Iv}{2\pi}\ln\frac{a+b}{a-b}>0}.
\]
\crosscheck 电动势对有向路径的符号与开路静电电势差的定义不同。$a>b$ 防止导体弧跨越无限长载流直线的奇点；结果只取决于端点横坐标，不依弧的具体形状。

\sourcetest{6.4.1}{166}{174}
\problemtext 证明：任意形状的定常载流导线段 $ab$ 置于空间均匀磁场中时，它所受总安培力等于把同电流沿端点间直线 $ab$ 放置所受的安培力。
\workings 将曲线按有向电流方向由点 $a$ 积到点 $b$，任一微小线元上
$d\bm F=I\,d\bm l\times\bm B$。因为 $\bm B$ 不随位置改变，可以提出积分号：
\[
\bm F=I\int_a^b d\bm l\times\bm B
=I\left(\int_a^b d\bm l\right)\times\bm B
=\boxed{I(\bm r_b-\bm r_a)\times\bm B}.
\]
而从 $a$ 到 $b$ 的直导线所受安培力恰为此式，证毕。
\crosscheck 此结论要求磁场\emph{空间均匀}且曲线各处电流相同；磁场不均匀时通常不能提出积分。闭合回路可取 $a=b$，故匀强场中总平动力为零，但力矩仍可能不为零。

\sourcetest{6.4.2}{166}{174}
\problemtext 半径为 $R$、磁导率近似真空的均匀圆柱长直导线中，总电流 $I$ 在截面上均匀分布。证明导线内部磁场能量每单位轴向长度为 $W_m=\mu_0 I^2/(16\pi)$。
\workings 半径 $r\le R$ 内包围电流 $I(r)=I r^2/R^2$，安培环路定律给出
\[
B(r)=\frac{\mu_0 I(r)}{2\pi r}
=\frac{\mu_0I r}{2\pi R^2}.
\]
真空或 $\mu_r\simeq1$ 时局部磁能密度 $w=B^2/(2\mu_0)$。单位长度圆柱内部的磁能为
\[
\begin{aligned}
W_m&=\int_0^R\frac{B(r)^2}{2\mu_0}(2\pi r)\,dr
=\frac{\mu_0 I^2}{4\pi R^4}\int_0^R r^3\,dr\\
&=\boxed{\frac{\mu_0 I^2}{16\pi}}.
\end{aligned}
\]

\crosscheck 积分仅包含\emph{导体内部}磁场能量；无限长导体外部磁场的径向能量积分若无外边界会出现对数发散。
% === END guide-selftest06-solutions.tex ===
